% Minimalt eksempel (MWE) på en øving med NTNUoving og ovingekstra.
%
% Denne filen gir oppgavesettet. Løsningsforslaget bygges fra mwe-lf.tex,
% som bare bytter til klassen NTNUlf og henter inn denne filen. Da står
% oppgaver og svar ett sted.
%
% Bygg med pdflatex to ganger (sidetall og hintliste trenger to runder):
%   pdflatex mwe.tex && pdflatex mwe.tex
%   pdflatex mwe-lf.tex && pdflatex mwe-lf.tex

\providecommand{\ovingklasse}{NTNUoving}
\documentclass[11pt,a4paper,norsk]{\ovingklasse}
\usepackage{ovingekstra}

% Valgfrie innstillinger, se README.md:
%\hinttittel{Hintliste}         % overskrift over hintene
%\hintkolonner                   % hint som tabell med «Hint 1» og «Hint 2»
%\hintliste                      % hint som nummerert liste, ett per oppgave
%\losningsetikett{Mulig svar:}   % tekst foran hvert svar
%\forfatterbredde{0.35}          % bredde på forfatterkolonnen i litteraturen

\fag{TBA4155 Byggeprosess: Tidligfasen i prosjekter}
\semester{Høst 2026}
\institutt{Institutt for bygg- og miljøteknikk}
\ovingnummer{1}
\ovingtittel{Eksempeløving}

\begin{document}

\ovingtittelblokk

\noindent\textbf{Litteratur:}

\begin{litteratur}
Ola Nordmann (2025). & ``Tittel på kapittel''. Forlag. \\ \midrule
Kari Nordmann og Per Hansen (2024). & ``Tittel på artikkel''. Tidsskrift 1 (2) s. 10-20. \\
\end{litteratur}

\bolk{Begrep og termer}

% \oppgave nummereres automatisk. \hint kan brukes én eller flere ganger
% per oppgave og samles i hintlisten bakerst. \losning vises bare i
% løsningsforslaget.
\oppgave{Forklar begrepet forventningsverdi.}
\hint{Kapittel 2.}
\hint{Middelverdi.}
\losning{Forventningsverdien er middelverdien, altså $x$-koordinaten til
tyngdepunktet for sannsynlighetsfordelingen.}

\oppgave{Finn standardavviket for summen av $\sigma_1 = 3$ og $\sigma_2 = 4$.}
\hint{Summer variansene.}
\losning{$\sigma = \sqrt{3^2 + 4^2} = 5$.}

\bolk{Anvendelse}

% Delspørsmål med a), b) inne i en oppgave, hver med sitt svar.
\oppgave{%
\begin{enumerate}[label=\alph*),leftmargin=*,topsep=0pt,itemsep=4pt]
  \item Hva er en styringsramme?
        \losning{Rammen som bevilges for å gjennomføre prosjektet, gjerne P50.}
  \item Hva er en kostnadsramme?
        \losning{Det beløpet som bør avsettes for å finansiere prosjektet,
        gjerne P85.}
\end{enumerate}}
\hint{Forelesning.}

% Navngitte delspørsmål med \sporsmaal.
\oppgave{Oppgaven er hentet fra eksamen. Svar kort.}
\hint{Se løsningsforslaget.}

\sporsmaal{a}
Hva menes med \textit{optimism bias}?
\losning{Å bedømme fremtidige hendelser i et mer positivt lys enn erfaring
tilsier.}

\sporsmaal{b}
Hva menes med \textit{strategic misrepresentation}?
\losning{Å undervurdere kostnader eller overvurdere nytte med overlegg.}

% Figurer settes inn som vanlig, for eksempel originalbilder fra kilden:
%\begin{center}
%\includegraphics[width=0.8\textwidth]{figurer/figur.png}
%\end{center}

\clearpage
\appendix
\hinttabell

\end{document}
